\section{Memory-preserving regularization and stability}
\label{sec:r79-label-restart}

Reachable supports can change at a parameter boundary, and the exact
potentials in Example~\ref{ex:r79-unbounded-reachable-witness} need
not be uniformly bounded.  We now regularize the controller without
adding a label or resetting the physical system.  The operation is
an occasional return to its already counted initial label.  The
all-initial-state cost quantifier makes this useful even when the
physical posterior and the fixed hidden model remain unknown.

The primitive fallback, centered effect estimate
\eqref{eq:effect-distance}, and compact-family completion in
Theorem~\ref{thm:compact-risk} are retained predecessors.  The
additional requirement here is that regularization and perturbation
use exactly the existing labels; no finite-history expansion is
allowed in the fixed-memory comparison.  Classical robust FSC
synthesis \cite{CubuktepeEtAl21,GaleslootEtAl25} and quantum cost
verification \cite{LiuZhouBartheYing25,AvanziniEtAl22} provide the
surrounding theory, but the bounds below are proved directly.

\subsection{Returning to the initial controller label}

\begin{theorem}[Regularization with the same live-label budget]
\label{thm:r79-label-restart}
Suppose an $M$-label controller $\gamma=(p,u)$ has robust cost at
most $b<\infty$ for the compact model family of
Section~\ref{sec:r79-quantum-frontier}.  For $0<\varepsilon\le1$
replace only its post-failure update rows by
\begin{equation}\label{eq:r79-label-restart}
 u^{k,a,y,\varepsilon}_\ell
       =(1-\varepsilon)u^{k,a,y}_\ell
                        +\varepsilon\mathbf1_{\{\ell=1\}}.
\end{equation}
The command rows and prescribed initial label are unchanged.  If
$\varepsilon b<c_*$, the resulting controller $\gamma_\varepsilon$
has robust initial cost at most
\begin{equation}\label{eq:r79-restart-initial-bound}
 U_\varepsilon=\frac{b}{1-\varepsilon b/c_*}.
\end{equation}
Moreover it has uniform cost at most
\begin{equation}\label{eq:r79-restart-global-bound}
 K_\varepsilon=\frac{c^*}{\varepsilon}+U_\varepsilon
\end{equation}
from every physical density matrix at every controller label, for
every fixed model.  Its exact cost operators can therefore be
chosen on the full product algebra and satisfy
\begin{equation}\label{eq:r79-global-witness}
 W^\theta=K^\theta_{\gamma_\varepsilon}
                +(\mathsf T^\theta_{\gamma_\varepsilon})^*(W^\theta),
 \qquad 0\preccurlyeq W^\theta_k\preccurlyeq K_\varepsilon I,
 \qquad W^\theta_1\preccurlyeq U_\varepsilon I.
\end{equation}
\end{theorem}
\begin{proof}
Realize \eqref{eq:r79-label-restart} by a fresh independent coin
after every failure: with probability $\varepsilon$ write label
$1$, and otherwise use the original update.  No physical operation
is performed by this coin.  The label update already belongs to
the declared controller interface.

Start from label $1$ with any physical state, and stop a segment
at success or at the first successful restart coin.  Up to that
coin its command and report law is the original controller's law
cut at an additional stopping time.  Its expected cost is at most
$b$.  The original controller has expected duration at most
$b/c_*$.  The probability of a restart before success is consequently
at most
\[
 r_\varepsilon=\varepsilon b/c_*.
\]
For example this follows by summing, over failure calls, the
probability of reaching that call without an earlier restart and
then obtaining a successful restart coin; deleting the restriction
on earlier restart coins only increases the sum to
$\varepsilon$ times the original expected number of failures.

After a restart the physical state can be arbitrary, but the
controller again has label $1$ and the same fixed model.  Conditional
on every such segment history, the all-initial-state and all-model
bound applies again.  Thus the probability of reaching segment
$j+1$ is at most $r_\varepsilon^j$, and its conditional expected
cost is at most $b$.  Summing finitely many segments and then using
monotone convergence gives $b\sum_{j\ge0}r_\varepsilon^j$.
This proves \eqref{eq:r79-restart-initial-bound} without assuming
the finiteness of the regularized cost in advance.

Now start at any other label and any physical state.  Before its
first success or restart, the probability of surviving $n$ calls
without either event is at most $(1-\varepsilon)^n$: each
nonterminal call offers an independent restart coin.  The expected
number of calls in that first segment is at most $1/\varepsilon$,
and its expected cost is at most $c^*/\varepsilon$.  Any following
segment starts at label $1$, whose total remaining cost is at most
$U_\varepsilon$.  This gives
\eqref{eq:r79-restart-global-bound}.

All blocks are now proper at finite cost.  The positive cost floor
and the order argument in Lemma~\ref{lem:r79-reachable-corner}
make the full product map transient.  Its geometric cost series
gives \eqref{eq:r79-global-witness}, and the cost bounds just
proved give the two operator inequalities.  The estimate is
uniform in the fixed model and does not require the observer to
know that model.
\end{proof}

The state of the physical experiment is unchanged at a controller
restart.  In particular, \eqref{eq:r79-label-restart} does not add
a freely available reset operation to a nonreset instrument.
Nor does it add a clock, a stored seed, or another live label.

\subsection{Perturbations of the primitive instrument}

To compare two primitive models $\theta,\vartheta$ with the same
physical and report spaces, include the scalar success outcome
and form the trace-preserving output map
\[
 \mathsf N^\theta_a(\rho)=
       \bigl((\Phi^\theta_{a,y}(\rho))_{y\in\mathcal Y},
                         \tr(B^\theta_a\rho)\bigr).
\]
The norm on this direct sum is the sum of the block trace norms.
Define
\begin{align}
 d_{\rm inst}(\theta,\vartheta)
  &=\max_a\sup_{\rho\succeq0,\ \tr\rho=1}
       \tfrac12\|\mathsf N^\theta_a(\rho)
                      -\mathsf N^\vartheta_a(\rho)\|_1,
                 \label{eq:r79-instrument-distance}\\
 d_{\rm cost}(\theta,\vartheta)
  &=\max_a\|G^\theta_a-G^\vartheta_a\|.
                 \label{eq:r79-cost-distance}
\end{align}
These distances compare the declared physical maps and mean costs.
They make no measurement available to the controller.

\begin{lemma}[Perturbation of a bounded cost witness]
\label{lem:r79-centered-perturbation}
Fix one counted controller and suppose
$0\preccurlyeq W_k\preccurlyeq KI$ at every live label.  If
$d_{\rm inst}(\theta,\vartheta)\le\eta$, then
\begin{equation}\label{eq:r79-adjoint-perturbation}
 -\eta KI\preccurlyeq
   ((\mathsf T^\vartheta_\gamma)^*
       -(\mathsf T^\theta_\gamma)^*)(W)_k
                  \preccurlyeq\eta KI
 \qquad(1\le k\le M).
\end{equation}
If in addition $d_{\rm cost}(\theta,\vartheta)\le\zeta$ and $W$
is a global supersolution for model $\theta$, then, provided
$\delta=\zeta+\eta K<c_*$,
\begin{equation}\label{eq:r79-scaled-witness}
 \frac{W}{1-\delta/c_*}
\end{equation}
is a global supersolution for model $\vartheta$, whose call costs
also satisfy the lower bound $c_*I$.
\end{lemma}
\begin{proof}
Fix a current label $k$ and an initial density matrix $\rho$.
After a sampled command $a$, regard the block observable on report
$y$ as $\sum_\ell u^{k,a,y}_\ell W_\ell$, and put the observable
zero on success.  Every block lies between zero and $KI$.
The two normalized output states have a difference $D$ of trace
zero.  Hence, for such an observable $Z$,
\[
 |\tr(ZD)|
   =|\tr((Z-\tfrac K2 I)D)|
   \le\tfrac K2\|D\|_1\le K\eta.
\]
Average over the same command row and then range over $\rho$.
The resulting scalar inequalities are exactly
\eqref{eq:r79-adjoint-perturbation}.  Including success in the
normalized output is what justifies the factor $1/2$; an
unnormalized failure-output difference need not have trace zero.

The source supersolution and the cost difference now give
\[
 W-(\mathsf T^\vartheta_\gamma)^*(W)
       \succeq K^\vartheta_\gamma-\delta I.
\]
Put $\alpha=(1-\delta/c_*)^{-1}$.  Since
$K^\vartheta_\gamma\succeq c_*I$ and
$(\alpha-1)c_* =\alpha\delta$, multiplying this inequality by
$\alpha$ gives
$\alpha W-(\mathsf T^\vartheta_\gamma)^*(\alpha W)
 \succeq K^\vartheta_\gamma$.
This proves \eqref{eq:r79-scaled-witness}.
\end{proof}

For compact families $\mathcal U,\mathcal V$ of primitive models,
write $V_M(\mathcal U)$ and $V_M(\mathcal V)$ for their respective
frontiers.  A directed comparison from $\mathcal U$ to $\mathcal V$
means that for every $\vartheta\in\mathcal V$ there is a
$\theta\in\mathcal U$ with both
$d_{\rm inst}(\theta,\vartheta)\le\eta$ and
$d_{\rm cost}(\theta,\vartheta)\le\zeta$.  The same paired models
must satisfy both bounds.

\begin{theorem}[Uniform stability of the complete memory frontier]
\label{thm:r79-frontier-stability}
Suppose the source family $\mathcal U$ is feasible and has common
cost bounds $c_*,c^*$.  Suppose the target family $\mathcal V$ has
the same lower cost bound and admits the directed comparison just
specified.  Put $b=V_M(\mathcal U)$.  For any
$0<\varepsilon\le1$ satisfying
\[
 \varepsilon b<c_*,\qquad
 \zeta+\eta K_\varepsilon<c_*,
\]
where $U_\varepsilon,K_\varepsilon$ are
\eqref{eq:r79-restart-initial-bound}--
\eqref{eq:r79-restart-global-bound}, one has
\begin{equation}\label{eq:r79-directed-stability}
 V_M(\mathcal V)
       \le\frac{U_\varepsilon}
                 {1-(\zeta+\eta K_\varepsilon)/c_*}.
\end{equation}
The target controller has the same $M$ labels as the source
optimizer.

In particular assume both families are feasible, have the common
cost bounds, have a one-row cost bound $B$, and admit the comparison
in both directions.  If $0<\eta\le1$ and
\begin{equation}\label{eq:r79-stability-smallness}
 \sqrt\eta\le\frac{c_*}{2B},\qquad
 \zeta+c^*\sqrt\eta+2B\eta\le\frac{c_*}{2},
\end{equation}
then, simultaneously for every $M\ge1$,
\begin{equation}\label{eq:r79-uniform-modulus}
 |V_M(\mathcal U)-V_M(\mathcal V)|
 \le\frac{10B^2+4Bc^*}{c_*}\sqrt\eta
                       +\frac{4B}{c_*}\zeta.
\end{equation}
For $\eta=0$ the same bound holds with the first term zero whenever
$0\le\zeta\le c_*/2$.
Consequently the entire family of finite-memory values is locally
uniformly continuous in the primitive model family, including at
changes of reachable support.
\end{theorem}
\begin{proof}
Take an attaining source controller, supplied by
Theorem~\ref{thm:r79-quantum-frontier}, and apply the label restart
of Theorem~\ref{thm:r79-label-restart}.  It is one common controller
for every source model and has global witnesses bounded by
$K_\varepsilon$.  For a target model $\vartheta$, choose its
paired source model $\theta$ only in the proof.  The preceding
lemma transfers that source witness to a target supersolution.
Its initial block is bounded by the right-hand side of
\eqref{eq:r79-directed-stability}.  The controller probabilities
do not depend on either member of this pair.  Applying the
certificate theorem separately to every fixed target model gives
one uniform robust bound and proves the directed inequality.
This is an offline comparison of complete fixed models; it does
not allow a model to change during the experiment or reveal a
model to the observer.

For the explicit two-sided estimate set $\varepsilon=\sqrt\eta$.
Since $b\le B$, the first smallness condition gives
\[
 U_\varepsilon\le2B,
 \qquad U_\varepsilon-b\le\frac{2B^2}{c_*}\sqrt\eta,
 \qquad
 \delta:=\zeta+\eta K_\varepsilon
       \le\zeta+c^*\sqrt\eta+2B\eta\le c_*/2.
\]
It follows that
\begin{align*}
 \frac{U_\varepsilon}{1-\delta/c_*}-b
 &\le\frac{2B^2}{c_*}\sqrt\eta
                       +\frac{4B}{c_*}\delta\\
 &\le\frac{2B^2+4Bc^*}{c_*}\sqrt\eta
          +\frac{8B^2}{c_*}\eta+\frac{4B}{c_*}\zeta.
\end{align*}
Use $\eta\le\sqrt\eta$ and then apply the same argument in the
opposite direction.  This proves
\eqref{eq:r79-uniform-modulus}.  For $\eta=0$, take
$\varepsilon\downarrow0$ in the directed inequality; the
resulting difference is at most
$b\zeta/(c_*-\zeta)\le2B\zeta/c_*$ in each direction, which
implies the stated bound.

Finally fix any feasible compact source family.  Its full-support
one-row maps in Theorem~\ref{thm:r79-dark-feasibility} have
invertible resolvents at every source model.  Continuity and
compactness give a common neighborhood in which the same row
remains proper with a bounded cost, and the positive cost bounds
persist locally within the class of valid instruments.  Thus a
common $B,c_*,c^*$ can be used for all sufficiently close compact
families.  The quantitative estimate proves the asserted local
uniform continuity.  No continuity of the optimizing supports,
labels, or individual unregularized potentials is required.
\end{proof}

The modulus is a proved upper estimate.  No optimality of its
square-root exponent is asserted.  Its constants do not grow with
the controller label budget; their dependence on the physical
experiment is displayed through the positive cost bounds and the
one-row cost cap.  The reset-specific continuity theorem remains
available with its stronger scalar structure.

\subsection{Relation to unrestricted visible history}

Let $V_{\rm hist}(\mathcal U)$ be the infimum of the robust cost
over all visible-history policies with model-independent causal
command kernels.  Success stopping and the physical observation
port are unchanged.  The following completion argument identifies
the limit of the counted-memory frontier; it does not assert
finite saturation.

\begin{corollary}[Finite-memory approximation to unrestricted control]
\label{cor:r79-history-limit}
For a feasible compact model family with one-row bound $C_0$,
\begin{equation}\label{eq:r79-history-limit}
 V_M\downarrow V_{\rm hist}\quad\hbox{as }M\to\infty.
\end{equation}
More explicitly, for every integer $T\ge1$ let
\[
 M_T=1+\sum_{t=0}^{T-1}(|\mathcal A||\mathcal Y|)^t.
\]
Then
\begin{equation}\label{eq:r79-history-rate}
 V_{\rm hist}\le V_{M_T}
       \le V_{\rm hist}
                    \left(1+\frac{C_0}{c_*T}\right).
\end{equation}
The initial label, every retained finite history, and the fallback
row are all included in $M_T$.  The uniform primitive stability
bound \eqref{eq:r79-uniform-modulus} also holds for
$V_{\rm hist}$.
\end{corollary}
\begin{proof}
Take any visible-history policy of robust cost $b<\infty$.  Follow
its causal kernels for at most $T$ calls, recording the preceding
command/report pairs in the counted labels.  On failure at the
$T$th call, switch to the common one-row policy of
Theorem~\ref{thm:r79-dark-feasibility}.  A rooted history tree has
at most $\sum_{t=0}^{T-1}(|\mathcal A||\mathcal Y|)^t$ nonterminal
prefixes; one additional label implements the fallback row.  Thus
this is an actual controller with at most $M_T$ labels and a pure
initial label.  The depth is stored in the history label rather
than supplied by a free clock.

The expected prefix cost is at most $b$.  The original policy has
expected duration at most $b/c_*$, so its probability of surviving
$T$ calls is at most $b/(c_*T)$.  On that event the actual physical
posterior is some density matrix and the model remains one fixed
member of the family.  The common all-initial-state fallback bound
therefore makes the additional cost at most
$C_0b/(c_*T)$.  The constructed controller has robust cost at most
$b(1+C_0/(c_*T))$.

Take the infimum over the original policies; an attaining
unrestricted policy is not assumed.  Every finite controller induces model-independent visible-history
command kernels.  Indeed, conditional on a fixed command/report
history, the physical branch likelihood is common to every compatible
private label path and cancels from its conditional distribution;
only the controller probabilities remain.  Arbitrary definitions on
zero-probability histories cause no change to a performed law.  Thus
finite controllers belong to the comparison class, yielding the
reverse inequality in \eqref{eq:r79-history-rate}.  Monotonicity in the label budget
and this approximating sequence prove
\eqref{eq:r79-history-limit}.  Finally pass to the infimum over
$M$ in the two-sided bound
\eqref{eq:r79-uniform-modulus}, whose constants are independent of
$M$.  This proves its asserted extension to $V_{\rm hist}$.
\end{proof}

The finite-history truncation and fallback are a completion device
within the inherited causal framework.  The new stability argument
uses a different operation: randomizing the existing label updates
to obtain global cost witnesses at the same fixed memory.  Neither
operation grants the controller a quantum reset, an unknown-state
oracle, or knowledge of the fixed hidden model.

\subsection{A uniform success floor: a Lipschitz frontier}

The general modulus allows commands which can fail with probability
one at some physical states.  A uniform lower success effect gives
a stronger estimate, directly from the classical resolvent argument.
The statement below records its exact dependence on the primitive
success margin and retains the same live-label budget throughout.

\begin{corollary}[Positive-success refinement]
\label{cor:r79-positive-success}
Let $\mathcal U,\mathcal V$ be compact model families on the same
interface and the positive-cost conventions of
Section~\ref{sec:r79-quantum-frontier}, with common ceiling $c^*$ and paired
primitive discrepancy bounds $\eta,\zeta$ in both directions.
Assume that for some $0<\kappa\le1$,
\begin{equation}\label{eq:r79-success-floor}
 B^\theta_a\succeq\kappa I
 \qquad\text{for every command and every model in }
                         \mathcal U\cup\mathcal V.
\end{equation}
Then, simultaneously for every $M\ge1$,
\begin{equation}\label{eq:r79-lipschitz-frontier}
 |V_M(\mathcal U)-V_M(\mathcal V)|
          \le\frac{\zeta}{\kappa}
                         +\frac{c^*\eta}{\kappa^2}.
\end{equation}
The same estimate holds for the unrestricted visible-history
values.  There is no small-error restriction in this inequality.

For either individual family, with its one-row fallback bound
$C_0\le c^*/\kappa$ and the counted $M_T$ of
Corollary~\ref{cor:r79-history-limit},
\begin{equation}\label{eq:r79-geometric-history}
 0\le V_{M_T}-V_{\rm hist}
                     \le C_0(1-\kappa)^T\qquad(T\ge1).
\end{equation}
This rate is exponential in the history horizon.  When
$|\mathcal A||\mathcal Y|>1$, the displayed history construction
uses exponentially many labels in that horizon, so the statement
does not give exponential convergence in the label count.

For $0<\kappa<1$, the universal local coefficients
$1/\kappa$ for cost perturbations and $c^*/\kappa^2$ for
instrument perturbations in \eqref{eq:r79-lipschitz-frontier}
cannot be reduced.  One-dimensional, one-command, one-model
experiments suffice to show this, at every live-label budget.
\end{corollary}
\begin{proof}
For any counted controller $\gamma$, normalization of its update
rows and \eqref{eq:r79-success-floor} give the block inequality
\[
 (\mathsf T^\theta_\gamma)^*(I)
       \preccurlyeq(1-\kappa)I.
\]
Positivity implies that this adjoint contracts the Hermitian block
operator norm by at most $1-\kappa$.  Thus every controller,
from every label and every physical state, is proper and has
exact cost potential
\[
 W^\theta=(I-(\mathsf T^\theta_\gamma)^*)^{-1}
                                     K^\theta_\gamma,
 \qquad 0\preccurlyeq W^\theta\preccurlyeq(c^*/\kappa)I.
\]
The Neumann inverse has norm at most $1/\kappa$ on Hermitian
blocks.  No label restart or additional controller state is needed.

For paired models $\theta,\vartheta$, subtract the two cost
equations to obtain
\begin{align*}
 W^\vartheta-W^\theta
  ={}&(I-(\mathsf T^\vartheta_\gamma)^*)^{-1}
       \bigl[K^\vartheta_\gamma-K^\theta_\gamma\\
   &\hspace{38mm}
       +((\mathsf T^\vartheta_\gamma)^*
                    -(\mathsf T^\theta_\gamma)^*)W^\theta\bigr].
\end{align*}
Lemma~\ref{lem:r79-centered-perturbation} bounds the bracket by
$\zeta+\eta c^*/\kappa$ in Hermitian operator norm.  Therefore
\[
 \|W^\vartheta-W^\theta\|
             \le\zeta/\kappa+c^*\eta/\kappa^2.
\]
This bounds the initial-block expected costs for every density
matrix.  Taking the worst model using the declared pairings, then
the infimum over the same controller class, proves one direction
of \eqref{eq:r79-lipschitz-frontier}; exchange the families for
the other direction.  The bound is independent of $M$.  Passing
to the limit in Corollary~\ref{cor:r79-history-limit} proves the
history assertion.

For the approximation rate, every visible-history policy has
conditional success probability at least $\kappa$ at each call,
irrespective of its history, unknown physical state, or fixed
model.  Its probability of surviving $T$ calls is at most
$(1-\kappa)^T$.  Truncate a finite-cost competing policy at that
horizon and execute the common fallback on its open branches.
The expected prefix cost is at most the original cost, and the
added cost is at most $C_0(1-\kappa)^T$.  The same history tree
and one fallback label give exactly the previous count $M_T$.
Taking the infimum over competing policies proves
\eqref{eq:r79-geometric-history}; an attaining unrestricted
policy is not required.

Finally take a scalar experiment with one command, one failure
report and mandatory success stopping.  Let the source success
probability be $\kappa+\eta$ and its deterministic call cost be
$C>0$, while the target success probability is $\kappa$ and its
call cost is $C+\zeta$, where
$0\le\eta\le1-\kappa$ and $\zeta\ge0$.
The instrument distance is exactly $\eta$, the cost distance is
exactly $\zeta$, and the common ceiling may be taken to be
$c^*=C+\zeta$.  No controller can affect these geometric waiting
times.  Their exact difference, at every memory, is
\begin{equation}\label{eq:r79-lipschitz-local-sharpness}
 \frac{C+\zeta}{\kappa}-\frac{C}{\kappa+\eta}
       =\frac{\zeta}{\kappa}
                     +\frac{C\eta}{\kappa(\kappa+\eta)}.
\end{equation}
Setting $\eta=0$ attains the cost coefficient exactly.  Setting
$\zeta=0$ and letting $\eta\downarrow0$ gives the ratio
$C/\kappa^2=c^*/\kappa^2$ for instrument error.  This proves
the two stated local sharpness assertions.
\end{proof}

The extra hypothesis \eqref{eq:r79-success-floor} is stronger than
absence of a common dark subspace: a jointly feasible detector
family may have a zero success effect on some state for each
individual command.  The corollary applies to the stated positive
success class only.  Its matched Lipschitz constants do not prove
sharpness of the square-root modulus in the general class, and
the resolvent method itself is not claimed as a new perturbation
principle.
